\subsection*{Geometric Intuition}

Riemann promotes Cauchy's tag-independence theorem into a definition. A
function is integrable when every sufficiently fine staircase, even with
adversarial tags, is forced near the same number.

\begin{remark*}[Engineering reading]
The sampling convention is no longer chosen by the author of the definition.
It is chosen by a universal quantifier. A Riemann-integrable function is one
whose accumulated area survives every possible choice of sample heights once
the slabs are fine enough.
\end{remark*}

\begin{definitionbox}{Definition (Riemann Sum)}
\begin{definition}[Riemann Sum]
\label{def:riemann-sum}
For a tagged partition \((P,\xi)\), define
\[
  S(f,P,\xi)=\sum_{i=1}^{n}f(\xi_i)\Delta x_i.
\]
\end{definition}
\end{definitionbox}

\begin{definitionbox}{Definition (Riemann Integral)}
\begin{definition}[Riemann Integral]
\label{def:riemann-integral}
A function \(f:[a,b]\to\mathbb R\) is \emph{Riemann-integrable} with value
\(L\) if
\[
  \forall\varepsilon>0\ \exists\delta>0\ \forall(P,\xi)\,
  \bigl(\|P\|<\delta\Rightarrow |S(f,P,\xi)-L|<\varepsilon\bigr).
\]
\end{definition}
\end{definitionbox}

\begin{remark*}[The one symbol that changed everything]
Cauchy quantifies over choppings. Riemann quantifies over choppings and every
possible choice of heights:
\[
  \forall P\qquad\leadsto\qquad \forall(P,\xi).
\]
That extra quantifier removes the arbitrary left endpoint and turns tag
independence into the price of admission.
\end{remark*}

\begin{remark*}[Payoff]
The definition is no longer tied to continuity. It can admit bounded functions
with discontinuities, because the question is not whether \(f\) is continuous
but whether all fine tagged sums are forced into one narrow numerical corridor.
The Cauchy criterion below makes that statement intrinsic: the value need not
be known in advance.
\end{remark*}

\begin{remark*}[The ruler function example]
Thomae's function, also called the ruler function, is
\[
  T(x)=
  \begin{cases}
    1/q,& x=p/q\text{ in lowest terms},\\
    0,& x\notin\mathbb Q.
  \end{cases}
\]
It is discontinuous at every rational and continuous at every irrational, so
its discontinuities are dense. Nevertheless its Riemann integral on any compact
interval is \(0\). The reason is that the rational points where \(T\) can be
large form only a finite set above any fixed height threshold. Fine partitions
can isolate those finitely many tall spikes, while all remaining slabs have
uniformly small height. This is exactly the kind of function Cauchy's
continuity proof could not diagnose, but Riemann's tag-independent definition
can admit.
\end{remark*}

\begin{theorembox}{Theorem (Continuous Implies Riemann-Integrable)}
\begin{theorem}[Continuous Functions Are Riemann-Integrable]
\label{thm:continuous-riemann-integrable}
Every continuous function \(f:[a,b]\to\mathbb R\) is Riemann-integrable.
\end{theorem}
\end{theorembox}

\begin{theorembox}{Theorem (Linearity)}
\begin{theorem}[Linearity of the Riemann Integral]
\label{thm:riemann-integral-linearity}
If \(f\) and \(g\) are Riemann-integrable on \([a,b]\), then
\[
  \int_a^b(\alpha f+\beta g)\,dx
  =
  \alpha\int_a^b f\,dx+\beta\int_a^b g\,dx .
\]
\end{theorem}
\end{theorembox}

\begin{theorembox}{Theorem (Monotonicity)}
\begin{theorem}[Monotonicity of the Riemann Integral]
\label{thm:riemann-integral-monotonicity}
If \(f\) and \(g\) are Riemann-integrable on \([a,b]\) and \(f(x)\le g(x)\)
for every \(x\in[a,b]\), then
\[
  \int_a^b f\,dx\le \int_a^b g\,dx .
\]
\end{theorem}
\end{theorembox}

\begin{theorembox}{Theorem (Triangle Inequality)}
\begin{theorem}[Triangle Inequality for the Riemann Integral]
\label{thm:riemann-integral-triangle-inequality}
If \(f\) is Riemann-integrable on \([a,b]\), then \(|f|\) is
Riemann-integrable and
\[
  \left|\int_a^b f\,dx\right|\le \int_a^b |f|\,dx .
\]
\end{theorem}
\end{theorembox}

\begin{theorembox}{Theorem (Interval Additivity)}
\begin{theorem}[Interval Additivity for the Riemann Integral]
\label{thm:riemann-integral-interval-additivity}
If \(c\in[a,b]\), then \(f\) is Riemann-integrable on \([a,b]\) if and only if
it is Riemann-integrable on \([a,c]\) and \([c,b]\), and in that case
\[
  \int_a^b f\,dx=\int_a^c f\,dx+\int_c^b f\,dx .
\]
\end{theorem}
\end{theorembox}

\begin{theorembox}{Theorem (Cauchy Criterion)}
\begin{theorem}[Cauchy Criterion for Riemann Integrability]
\label{thm:riemann-cauchy-criterion}
A bounded function is Riemann-integrable if and only if for every
\(\varepsilon>0\) there is \(\delta>0\) such that any two tagged partitions
\((P,\xi)\) and \((Q,\eta)\) with \(\|P\|,\|Q\|<\delta\) have
\[
  |S(f,P,\xi)-S(f,Q,\eta)|<\varepsilon.
\]
\end{theorem}
\end{theorembox}

\begin{remark*}[Bug report]
Riemann fixed the arbitrary tag, but the definition is operationally hostile.
To verify integrability directly one must control all tagged partitions at
once. The patch request is to remove tags from the verification problem:
replace every possible sample height by the slab infimum and supremum, then
squeeze.
\end{remark*}

\begin{remark*}[Why this is hard to use]
Riemann is logically clean but not mechanically friendly. Refining a partition
does not make an arbitrary Riemann sum move monotonically, because every
refinement lets the tags move again. There is no single upper or lower
quantity improving under refinement. Darboux supplies exactly that missing
monotone control.
\end{remark*}
